\documentclass[a4paper, 11pt]{article}

\usepackage[margin=3cm]{geometry} 
\usepackage{amsthm, amsmath, amssymb}
\usepackage{array}
\usepackage{bussproofs}
\usepackage{booktabs}

\theoremstyle{definition}
\newtheorem{definition}{Definition}[section]
\newtheorem{example}[definition]{Example}
\newtheorem{theorem}[definition]{Theorem}
\newtheorem{lemma}[definition]{Lemma}
\newtheorem{prop}[definition]{Proposition}
\newtheorem{corollary}[definition]{Corollary}

% change spacing
\let\olditemize\itemize
\renewcommand{\itemize}{\olditemize\setlength{\itemsep}{0pt}\setlength{\parsep}{0pt}}
\let\oldenumerate\enumerate
\renewcommand{\enumerate}{\oldenumerate\setlength{\itemsep}{0pt}\setlength{\parsep}{0pt}}

% hilbert proofs
\newcounter{dline}
\newenvironment{deduction}[1][\textit{no hypotheses}]
    {\setcounter{dline}{0}%
     \[\begin{array}{r|l@{\qquad}l}
     & #1 \\ \cline{2-2}}
    {\end{array}\]}
\newcommand{\dl}[2]{\stepcounter{dline}\thedline & #1 & \text{#2} \\}

% tables
\newenvironment{comparison}[4][p{0.22\textwidth} p{0.36\textwidth} p{0.36\textwidth}]
    {\setcounter{factrow}{0}%
     \begin{center}
     \begin{tabular}{#1}
     \textsc{#2} & \textsc{#3} & \textsc{#4} \\[4pt]}
    {\end{tabular}
     \end{center}}


\begin{document}


\section{Number Systems}

\subsection{Notation}
When programming, we cannot use subscripts like ${17}_8$ to denote the base.
\begin{itemize}
    \item Binary denoted as \texttt{0b}, e.g. \texttt{0b110000011111}
    \item Octal (base 8) denoted as \texttt{0}, e.g. \texttt{06037}
    \item Decimal is without prefix, e..g \texttt{3103}
    \item Hex denoted as \texttt{0x}, e.g. \texttt{0xC1F}
\end{itemize}

\subsection{Representing negative numbers}
Two's complement: ``The two's complement of an N-bit number is defined as its complement with respect to $2^N$''
\[
\overset{\substack{2^3 \\ =}}{1000_2} - \overset{\substack{3 \\ =}}{011_2}
\;=\; 101 \quad (\text{represents } -3)
\]
\[
x = (x_{N-1}\, x_{N-2} \ldots x_0)_2
  = \left( -x_{N-1}\, 2^{N-1} + \sum_{i=0}^{N-2} x_i\, 2^i \right)
\]
\begin{center}
\begin{tabular}{rc}
    \toprule
    \textbf{Decimal} & \textbf{Two's comp.} \\
    \midrule
    $-4$ & 100 \\
    $-3$ & 101 \\
    $-2$ & 110 \\
    $-1$ & 111 \\
    $0$  & 000 \\
    $1$  & 001 \\
    $2$  & 010 \\
    $3$  & 011 \\
    \bottomrule
\end{tabular}
\end{center}
Why 2's-complement? Single zero reprsentation simplifies equality checking, and arithmetic is overall easier.

\subsection{Integers}
\begin{itemize}
    \item \textit{Signed} $n$-bit integers can represent $-2^{n-1}, \dots, 2^{n-1} - 1$
    \item \textit{Unsigned} $n$-bit integer can represent $0, \dots, 2^{n-1}$
\end{itemize}

\subsection{Checking overflows in $2$'s complement}
\begin{itemize}
    \item (pos) $+$ (neg) cannot overflow
    \item (pos) $-$ (pos) cannot overflow 
    \item (neg) $+$ (neg) should be negative. Carry bit turns it positive.
    \item (pos) $+$ (pos) should be positive. Carry bit turns it negative.
\end{itemize}

\subsection{Extending numbers}
Consider the 8-bit number $01001011_2$. To extend it to 16 bits, we pad the front with 0s: $0000000001001011_2$.
Now the number $11001011_2$: If unsigned, we pad with 0s as prior. If signed, we pad with 1s instead: $1111111111001011_2$.



\end{document}